\section{Reasoning-Augmented Extraction}

\texttt{Reasoned-RAG} adds a temporary reasoning interface to the shared extraction pipeline. Some GenSIE fields require local evidence assessment before a value can be produced: the model may need to identify relevant fragments, decide whether evidence is absent, normalize a phrase, or map cues to an enum label. Reasoning-augmented extraction gives each top-level field a structured place for that intermediate decision while preserving the final JSON interface.

In \texttt{Reasoned-RAG}, each top-level field of the target object is temporarily wrapped as an object with two properties:

\begin{verbatim}
{
  "reasoning": "...",
  "value": ...
}
\end{verbatim}

The \texttt{reasoning} property is a string. The \texttt{value} property keeps the original schema of the wrapped field. If the original field is an enum, the \texttt{value} must be one of the enum labels; if it is an array or object, the nested structure remains there; if it is nullable, the \texttt{value} may be JSON \texttt{null} when the schema permits it. The wrapper changes the temporary generation shape, not the semantic target.

The submitted reasoning mode is top-level only. It wraps root object properties, not every nested property recursively. If a top-level field is itself an object or an array of objects, the nested content remains inside \texttt{value} according to the original field schema. This keeps the prompt and response format manageable while asking the model to deliberate at the level of the main requested slots.

The prompt-side effect is a reasoned Pydantic-style schema view. Rather than seeing a top-level field simply as a value of type \(T\), the model sees it as a reasoned value, using the class \texttt{Reasoned[T]}. The extraction instructions ask the model to state what the field requests, identify supporting evidence or absence, and then commit to the final value. This is intended for fields involving nullability, enum classification, boolean inference, normalization, and arrays whose correct value may be empty.

The generation-side effect is equally important. The reasoning interface is part of the strict temporary response schema, so each wrapped top-level field must contain both \texttt{reasoning} and \texttt{value}. This prevents free-form reasoning from leaking outside the JSON object and keeps the final value constrained by the original field type. The wrapper adds a field-local explanation channel without relaxing the schema constraints on the answer.

After generation, postprocessing removes the wrapper. Each top-level object of the form \texttt{\{reasoning, value\}} is replaced by its \texttt{value}; reasoning strings are discarded before the output is returned or passed to self-consistency aggregation. If the expected wrapper structure is malformed and cannot be reliably unwrapped, the extraction record is treated as invalid rather than repaired semantically.

Reasoning-augmented extraction also affects RAG rendering. Examples retrieved for \texttt{Reasoned-RAG} are shown in the same temporary output format expected from the current extractor. A direct example with final field values would be inconsistent with a prompt that asks for top-level \texttt{\{reasoning, value\}} objects. Reasoning-mode examples therefore demonstrate both the evidence protocol and the placement of the final answer in the \texttt{value} slot, while still remaining demonstrations rather than evidence for the current instance.